% Concept-paper close: limitations + conclusion. CPU-only. Direct prose.
\section{Limitations}
\label{sec:c-close-limitations}
\begin{enumerate}
\item Single model, single seed: one reader, no family or seed sweep; generality untested.
\item Context capped at 16k; behaviour beyond unmeasured.
\item Text-only: verbatim only, no live key-value cache.
\item Inert pointers (M7): no measured retrieval effect; SemPointer selects the same blocks as dense retrieval.
\item No pointer-length sweep: fixed length/depth, capacity sensitivity unknown.
\item Small pilots: key-value n=10, wide uncertainty; multi-hop n=30, neither significant.
\item Single-8GB CPU-only rig: rig-specific latency/means, no generality claim.
\item Double-blind note: author-visible repository, deanonymizing; confirm model with chairs.
\end{enumerate}

\section{Conclusion}
\label{sec:c-close-conclusion}
SemPointer as assembled works: dense per-query routing answers at 10--15x fewer tokens than full context (deltas -0.37, -0.033, -0.012). Addressing distinct from retrieve-then-read remains unproven: no selection effect, undecided n=10 pilot, negative n=30 probe.

Verdict flips only on P1--P3: beating dense retrieval at matched budgets with a non-zero pointer effect, matching full context within uncertainty, or a key-value end-to-end reversing the deficit.

Future: larger multi-seed suites, pointer-length sweep, measured key-value end-to-end.
